\sectionquestion{What do we need from classes?}
\section[Refresher]{Refresher of object oriented programming}

\begin{frame}[fragile]{Classes}

    \inputminted{python}{code/refresher/small_refresher_of_classes_and_objects.py}

\end{frame}




\begin{frame}[fragile]{Inheritance: \textit{is a}}

    \inputminted{python}{code/refresher/inheritance_is_a.py}

\end{frame}




\begin{frame}[fragile]{Interfaces: \textit{can do}}

    \begin{columns}[T]
        \column{0.48\textwidth}
        \inputminted{python}{code/refresher/interface_contract.py}

        \column{0.48\textwidth}
        \inputminted{python}{code/refresher/interface_usage.py}
        \pause
        \inputminted{python}{code/refresher/interface_broken.py}
    \end{columns}

    \vfill

    \takeaway{An interface is a contract: which methods exist, not how they work.}
    \begin{center}
    Code that uses \texttt{Force} doesn't care which force it gets.
    \end{center}

    \vfill
\end{frame}




\begin{frame}[fragile]{Inheritance vs Interface}

    Python has no \texttt{interface} keyword. The same syntax means two things:

    \vfill
    \inputminted{python}{code/refresher/same_syntax.py}
    \vfill

    \pause

    \texttt{Force} is an interface because it holds \textbf{only promises}:
    \begin{itemize}
        \item Abstract methods: signatures, no code
        \item No attributes: no data
        \item Any number of classes can sign the same contract
    \end{itemize}

    \takeaway{Inheritance shares code. An interface only promises behaviour.}
    \begin{center}
    \small In this talk, every \texttt{ABC} is used as an interface.
    \end{center}
\end{frame}
